\scanpage{21}

d'ouverts $W_i$ \deleted{de} $X$ recouvrant $V_{(x)}$ \inserted{(donc en vertu du lemme on peut supposer qu'ils recouvrent $V$)} et des $f_i\in\Gamma(W_i,F)$, \unsure{telles|tels} que $y\in W_{ix}\deleted{\unsure{V_{(x)}|V_x}}\Rightarrow f_{iy}=f_y$.

Soit \deleted{$\unsure{f}$} $W'_{ij}$ l'ens des $y\in W_i\cap W_j=W_{ij}$ tels que $f_{iy}=f_{jy}$. Cet ens \deleted{est} \inserted{est ouvert et} contient $W_{ij}\cap V_{(x)}=W_{ij\,x}$

Par le lemme \deleted{\unsure{préc.}} déjà cité, \unsure{rapetissant} $X$ \deleted{$W\cap V$} autour de \deleted{$\unsure{(x)}$}$x$, on peut supposer que $W'_{ij}=W_{ij}$. Alors on dispose d'une vraie section. OK ....

% slanted
\margin{Généraliser \deleted{\unsure{ces}} les lemmes à des procédés de localisation plus généraux.....}

Prouvons maintenant $(\mathrm{i})\Rightarrow(\mathrm{ii})$ : En effet, \unsure{soit|soit le} $x\in Y$, \deleted{\unsure{considérons}} notons que par le \deleted{\unsure{coroll}} corollaire 4, $F$ satisfait à $P_1$ $(P_2)$ rel. à $\overline x$, donc \inserted{posant $V=X-\overline x$,} on a pour tt voisinage ouvert $W$ de $x$

\[
\Gamma(W,F)\longrightarrow\Gamma(W\cap V,F)
\]

est inj. (bij), donc

\[
\underset{W}{\unsure{\underline{\lim}|\varinjlim}}\ \Gamma(W,F)\mathrel{\unsure{-|\to}}\Gamma(W\cap V,F)
\]

est injectif (bij) or cet hom s'identifie à

\[
\Gamma(X_{(x)},F_{(x)})\longrightarrow\Gamma(X'_{(x)},F_{(x)})
\]

en appliquant deux fois le lemme 5, OK.

% archive sheet number: 20
